\chapter{Continuous-Time PIMC (\code{ctpimc})}
\label{ch:ctpimc}

The fourth engine removes the Trotter discretisation altogether: it samples
the path integral directly in \emph{continuous} imaginary time, representing
each spin as a piecewise-constant \emph{worldline} on the circle
$[0,\beta)$. There is no Trotter error, and the cluster updates it employs
mix particularly well in the strongly quantum (large-$\Gamma$) regime. The
implementation builds on D-Wave's \code{localPIMC} code (Apache-2.0,
credited in \code{NOTICE}).

\section{From $M\to\infty$ to worldlines}

Recall the discrete action Eq.~\eqref{eq:action}. Take
$M\to\infty$ at fixed $\beta$: the slice index becomes a continuous
imaginary time $\tau\in[0,\beta)$ and each site's spin becomes a function
$s_i(\tau)\in\{-1,+1\}$, periodic in $\tau$, changing value at a finite set
of \emph{kinks} (domain walls in time). In this limit:

\begin{itemize}
\item The classical term becomes an integral,
$\frac{\beta}{M}\sum_t\Ecl(s^t)\;\longrightarrow\;\int_0^\beta
\Ecl\bigl(s(\tau)\bigr)\,d\tau$.
\item The Trotter bonds become a \emph{kink fugacity}: each kink on a
worldline carries a weight $\propto\Gamma\,d\tau$, i.e.\ in the absence of
couplings the kinks of site $i$ form a \emph{Poisson process} of rate
$\Gamma$ on the time circle.
\end{itemize}

\begin{keyeq}[Continuous-time path-integral weight]
\begin{equation}
Z=\sum_{\text{worldline configs}}
\;\Gamma^{\,K}\,
\exp\!\left[-\int_0^\beta \Ecl\bigl(s(\tau)\bigr)\,d\tau\right],
\end{equation}
where $K$ is the total number of kinks and the ``sum'' is over all
piecewise-constant periodic $\{s_i(\tau)\}$.
\end{keyeq}

Large $\Gamma$ favours many kinks (strong quantum fluctuation, rapidly
flapping worldlines); $\Gamma\to0$ suppresses kinks and worldlines freeze
into constants --- the classical configurations.

\section{Cluster updates on worldlines}

The engine updates worldlines with Swendsen--Wang-style cluster moves over
worldline \emph{segments} (the intervals between kinks): segments of
neighbouring sites are stochastically bound together according to the
couplings acting during their overlap in time, and bound clusters are
flipped as a unit. Kinks are created and removed in the process, so both the
spin values and the kink structure are resampled. This is the
continuous-time analogue of the discrete SW move of
Section~\ref{sec:sw-moves}, without any discretisation scale.

Two knobs shape the proposals:
\begin{itemize}
\item \code{qubits\_per\_update} --- how many sites participate in one
cluster proposal (1 = single-site updates);
\item \code{qubits\_per\_chain} --- chain length for the multi-site chain
updates used by the lattice mode.
\end{itemize}

\section{General mode}

For arbitrary \code{DenseIsing}/\code{SparseIsing} problems the annealer
follows a standard $(\beta,\Gamma)$ schedule. At each schedule step the
couplings are rescaled to the dimensionless combinations the sampler
consumes,
\begin{equation}
\beta J_{ij},\qquad \beta h_i,\qquad \beta\Gamma ,
\end{equation}
then \code{sweeps\_per\_beta} cluster sweeps are run and the classical
projection of the current worldlines is recorded. Reported energies are
converted back to the \emph{original} Ising units (not scaled by $\beta$).

\begin{pycode}
import numpy as np
from qanneal import DenseIsing, SQASchedule, CTPIMCAnnealer

h = np.array([0.2, -0.3, 0.1])
J = np.array([[0.0, -1.0, 0.2],
              [-1.0, 0.0, 0.5],
              [ 0.2, 0.5, 0.0]])
ising = DenseIsing(h, J)

schedule = SQASchedule.from_vectors(
    betas  = np.linspace(0.5, 4.0, 40).tolist(),
    gammas = np.linspace(2.0, 0.05, 40).tolist(),
)
annealer = CTPIMCAnnealer(ising, schedule,
                          qubits_per_update=1, qubits_per_chain=1)
res = annealer.run(sweeps_per_beta=50, reads=4)
print(res.best_energy)
\end{pycode}

Through the high-level interface:

\begin{pycode}
result = solve(ising, method="ctpimc", reads=8,
               ctpimc_qubits_per_update=1)
\end{pycode}

\begin{warnbox}[$\Gamma$ must stay positive]
The continuous-time weight is singular at $\Gamma=0$ (zero kink rate: the
sampler cannot move). End the schedule at a small positive value
(e.g.\ $10^{-3}$), never exactly zero. When building tuned schedules for
CT-PIMC, pass \code{gamma\_end\_scale=0.04} to
\code{auto\_schedule\_sqa\_tuned} so the final $\Gamma$ is near-zero for a
clean classical projection, yet finite.
\end{warnbox}

\section{Lattice mode}

A second constructor reproduces the cylindrical-lattice experiments the
underlying sampler was designed for (triangular and square-octagonal
geometries):

\begin{pycode}
from qanneal import CTPIMCAnnealer

annealer = CTPIMCAnnealer(
    Lperiodic=12,          # lattice period (multiple of 6)
    inv_temp_over_J=2.0,   # beta / J
    gamma_over_J=0.6,      # Gamma / J
    initial_condition=0,   # 0 random, 1 ferro, -1 antiferro
    qubits_per_update=1,
    qubits_per_chain=1,
)
res = annealer.run(sweeps_per_beta=32768)
\end{pycode}

Here temperatures and fields are expressed in units of the lattice coupling
$J$, and reported energies are likewise in units of $J$.

\section{CT-PIMC versus SQA}

\begin{center}
\small
\begin{tabular}{@{}lll@{}}
\toprule
 & \textbf{SQA (Chapter~\ref{ch:sqa})} & \textbf{CT-PIMC}\\
\midrule
Imaginary time & $M$ discrete slices & continuous\\
Systematic error & $O(\beta^2\Gamma^2/M)$ Trotter error & none\\
Elementary move & spin / worldline / SW segment & worldline-segment clusters\\
High-$\Gamma$ mixing & needs large $M$ & natural regime\\
Replica / PT structure & yes (\code{sqapt}) & no\\
Sweep-level observers & yes & no\\
\bottomrule
\end{tabular}
\end{center}

Use CT-PIMC when you want discretisation-free sampling closest in spirit to
physical-annealer statistics; use SQA/SQAPT when you need the replica
machinery, adaptive $\Jp$ schedules, or detailed observer traces. Neither
dominates the other on all problems.
